\documentclass[10pt]{article}
\usepackage[a4paper,margin=2.1cm]{geometry}
\usepackage[T1]{fontenc}
\usepackage[utf8]{inputenc}
\usepackage{lmodern}
\usepackage{microtype}
\usepackage{xcolor}
\usepackage{enumitem}
\usepackage{tabularx}
\usepackage{longtable}
\usepackage{array}
\usepackage{amssymb}
\usepackage{hyperref}
\usepackage{fancyhdr}
\definecolor{tudblue}{RGB}{0,90,140}
\definecolor{draftred}{RGB}{165,35,35}
\definecolor{lightblue}{RGB}{232,242,247}
\hypersetup{colorlinks=true,linkcolor=tudblue,urlcolor=tudblue}
\setlength{\parindent}{0pt}
\setlength{\parskip}{0.55em}
\setlist{itemsep=0.25em,topsep=0.35em}
\renewcommand{\arraystretch}{1.22}
\pagestyle{fancy}
\fancyhf{}
\lhead{TU Darmstadt Ethics Commission}
\rhead{Draft for supervisor review}
\cfoot{\thepage}
\newcommand{\checked}{\ensuremath{\boxtimes}}
\newcommand{\unchecked}{\ensuremath{\square}}
\newcommand{\confirm}[1]{\textcolor{draftred}{\textbf{TO CONFIRM:} #1}}
\newcommand{\field}[2]{\textbf{#1}\par #2\par}
\newcommand{\draftbanner}{%
  \begin{center}
  \fcolorbox{draftred}{white}{\parbox{0.91\textwidth}{\centering
  \textcolor{draftred}{\textbf{DRAFT -- NOT FOR SUBMISSION OR DATA COLLECTION}}\\
  Administrative and data-management fields marked ``TO CONFIRM'' require supervisor approval.}}
  \end{center}}
\newcommand{\sectionbox}[1]{%
  \vspace{0.4em}\noindent\colorbox{lightblue}{\parbox{\dimexpr\linewidth-2\fboxsep\relax}{\textbf{#1}}}\vspace{0.35em}}

\lhead{Local Trait Lab}
\rhead{Supervisor meeting notes}
\begin{document}

\begin{center}
{\Large\bfseries Ethics Application Meeting Notes}\par
\vspace{0.25em}
{\large Meeting on 13 August 2026}\par
\end{center}

\sectionbox{Decisions needed from Simon / principal investigator}
\begin{enumerate}
  \item Confirm the principal investigator who will submit and sign the application.
  \item Confirm recruitment for approximately 100 adults: personal network, TU mailing lists, participant pool, social media, or a combination.
  \item Confirm compensation, or explicitly state that participation is unpaid. Avoid dependency relationships in which participation could feel obligatory.
  \item Confirm whether sessions are supervised, remote on participants' own Android phones, or both.
  \item Confirm the duration after a pilot. The ethics draft estimates 25--40 minutes plus possible model-download time; the app currently says about 15 minutes.
  \item Confirm the approved TU/PEASEC server, hosting location, encrypted upload, access-control list, backups, and responsible data handler.
  \item Confirm deletion periods for pseudonymized study data, anonymized data, and consent records.
  \item Confirm the last point at which participants can request deletion using their study ID.
  \item Confirm whether pseudonymized record-level data may be reused. Recommended default: no reuse beyond this project without renewed review.
  \item Confirm Dr. Nina Gerber's role: participating researcher, co-supervisor, principal investigator, or consultant.
\end{enumerate}

\sectionbox{Required app and data-flow fixes before data collection}
\begin{enumerate}
  \item Replace original filenames and local URIs in exports with random per-file identifiers. Filenames may contain names, account numbers, courses, or other personal information.
  \item Avoid direct production-study downloads from Hugging Face unless third-party IP-address disclosure is documented and approved. Prefer TU hosting, supervised preinstallation, or another approved route.
  \item Explicitly delete each temporary camera file immediately after classification and test this behavior. The current camera API creates a temporary image file.
  \item Decide whether optional reaction tracking covers the whole study. It currently ends when the pre-questionnaire page is left and therefore does not measure reactions to file inferences or the profile.
  \item Generate a random study ID, include it in the export, and give it to the participant for withdrawal requests.
  \item Let participants retain the complete information and privacy statement, not only the short in-app consent card.
  \item Require active answers or explicitly represent missing questionnaire responses rather than silently accepting visual slider defaults.
  \item Add ``prefer not to say'' for age, self-description for gender, and ``other'' for education.
  \item Verify that selected files, PDF previews, camera frames, and image bytes never enter exported JSON or server logs.
  \item Test encrypted transmission, certificate validation, failed uploads, local cleanup, and server deletion before recruitment.
\end{enumerate}

\sectionbox{Methodological questions}
\begin{enumerate}
  \item Keep IUIPC-8 and ATI as validated baseline scales; confirm scoring and analysis with Nina.
  \item Decide whether the two additional baseline items overlapping ATI or privacy concern are needed.
  \item Define whether optional emotion labels have a hypothesis or remain exploratory. They are not ground truth for participants' emotions.
  \item Decide whether approximately 20 files is mandatory, a target, or a maximum, and define handling of incomplete sessions.
  \item Confirm whether first-page-only PDF analysis is scientifically sufficient and disclose it consistently.
  \item Predefine primary outcomes, likely per-inference unexpectedness, sensitivity and perceived accuracy, plus final expectation, concern, reassurance, and access-willingness measures.
  \item Decide how malformed, failed, or clearly inaccurate Gemma outputs enter the analysis and whether participants can flag nonsensical results.
  \item Decide whether prompts avoid special-category inferences such as health, religion, political opinions, or sexual orientation, or whether these are central and require explicit consent and stronger safeguards.
\end{enumerate}

\sectionbox{Submission checklist}
\begin{enumerate}
  \item Do not collect data before a positive ethics vote.
  \item Submit one PDF ordered as: application form, participant information and consent, questionnaires, then other supporting material.
  \item Obtain the principal investigator's signature.
  \item Remove every red ``TO CONFIRM'' field and all internal draft notes.
  \item Ensure the app, consent, questionnaire appendix, and actual data flow agree exactly before recruitment.
\end{enumerate}

\end{document}
